\chapter{Inventario de Bienes y Herramientas}

El inventario al 31 de octubre de 2025 se verificó en conjunto con el operador y la presidencia saliente. El detalle completo se conserva en el archivo digital \emph{anexos/C\_listados/inventario\_operativo.xlsx} y en la carpeta física entregada, por lo que aquí se presentan los puntos esenciales para la entrega-recepción.

\section{Bienes durables}

Los bienes durables incluyen equipos, herramientas y mobiliario con vida útil mayor a un año. Se revisó su estado operativo, la ubicación y la persona responsable inmediata.

\subsection*{Puntos de control}
\begin{itemize}[noitemsep, topsep=0pt]
  \item Verificar que cada bien cuente con etiqueta o código de inventario y que coincida con el registro digital.
  \item Documentar observaciones sobre mantenimiento pendiente o reposición y dejarlas firmadas en el acta interna.
\end{itemize}
\textit{Nota: El registro completo por serie y condición se adjunta en el anexo indicado.}

\subsection*{Ejemplos de bienes verificados}
\begin{table}[H]
\centering
\caption{Bienes representativos (detalle completo en anexos)}
\label{tab:bienes_representativos}
\begin{tabular}{p{4.2cm} p{3cm} p{3cm} p{3cm}}
\toprule
\textbf{Descripción} & \textbf{Ubicación} & \textbf{Responsable} & \textbf{Estado} \\
\midrule
Máquina cloradora \textit{[Marca]} & Planta de tratamiento & Operador: [Nombre] & Operativa \\
Archivador metálico 4 gavetas & Oficina administrativa & Secretaria: [Nombre] & Bueno, requiere llave de repuesto \\
\bottomrule
\end{tabular}
\end{table}
\textit{El inventario completo de bienes durables, incluyendo fotografías y números de serie, permanece en \emph{anexos/C\_listados/}.}

\section{Insumos y stock operativo}

El stock operativo contempla productos químicos, tubería, accesorios y materiales de emergencia. Se recomienda actualizar el conteo inmediatamente después de la entrega para asegurar reposiciones oportunas.

\subsection*{Alertas y acciones inmediatas}
\begin{itemize}[noitemsep, topsep=0pt]
  \item Confirmar existencias de hipoclorito de calcio y programar compra cuando el saldo baje del mínimo operativo (25 kg).
  \item Revisar tubería y accesorios priorizando diámetros de 1'' y 3/4'' que se emplean en reparaciones de acometidas.
\end{itemize}
\textit{El balance detallado por código y proveedor se encuentra en la hoja de cálculo de anexos.}

\subsection*{Ejemplos de insumos críticos}
\begin{table}[H]
\centering
\caption{Insumos representativos (detalle completo en anexos)}
\label{tab:insumos_representativos}
\begin{tabular}{p{4.2cm} p{2.6cm} p{3cm} p{3.2cm}}
\toprule
\textbf{Insumo} & \textbf{Unidad} & \textbf{Cantidad física} & \textbf{Mínimo sugerido} \\
\midrule
Hipoclorito de calcio 65\% & kg & [Cantidad verificada] & 50 kg \\
Tubo PVC presión 1'' (6 m) & unidad & [Cantidad verificada] & 5 unidades \\
\bottomrule
\end{tabular}
\end{table}
\textit{Los saldos de otros insumos (válvulas, abrazaderas, adhesivos y materiales de construcción) constan en el anexo digital y se entregan en la bodega inventariada.}

\section{Seguimiento recomendado}

\begin{itemize}[noitemsep, topsep=0pt]
  \item Levantar inventario físico trimestral con participación del operador y del área administrativa para actualizar diferencias.
  \item Registrar entradas y salidas en formato firmado (libreta o sistema digital) y respaldar la información en Google Drive.
\end{itemize}
\textit{Cualquier novedad posterior deberá documentarse en acta complementaria y notificarse a la directiva entrante.}
